% !TeX program = xelatex
% Overleaf 编译器请选择 XeLaTeX。
\documentclass[a4paper,10pt]{article}

\usepackage[margin=1.25cm]{geometry}
\usepackage[UTF8,fontset=fandol]{ctex}
\usepackage{enumitem}
\usepackage{hyperref}
\usepackage{xcolor}
\usepackage{titlesec}
\usepackage{tabularx}
\usepackage{array}
\usepackage{microtype}
\usepackage{graphicx}

\definecolor{main}{HTML}{008C8C}
\definecolor{textgray}{HTML}{333333}
\definecolor{muted}{HTML}{666666}

\hypersetup{
    colorlinks=true,
    urlcolor=main,
    linkcolor=main
}

\pagestyle{empty}
\color{textgray}

\setlength{\parindent}{0pt}
\setlength{\parskip}{0pt}
\setlength{\tabcolsep}{0pt}
\renewcommand{\baselinestretch}{1.08}
\emergencystretch=2em

\setlist[itemize]{
    leftmargin=1.25em,
    itemsep=2.0pt,
    topsep=2.5pt,
    parsep=0pt,
    partopsep=0pt,
    label=\textcolor{main}{\small\textbullet}
}

\titleformat{\section}
{\large\bfseries\color{main}}
{}{0pt}{}
[\vspace{-1pt}\titlerule]

\titlespacing*{\section}{0pt}{5pt}{4pt}

\newcommand{\resumeItem}[1]{\item #1}

\newcommand{\skill}[2]{%
    \textbf{#1：}#2\par\vspace{1.5pt}
}

% #1 项目名
% #2 时间
% #3 技术栈
% #4 GitHub 完整链接
\newcommand{\projectTitle}[4]{%
    \textbf{\large #1}
    \hfill
    \textcolor{muted}{#2}\\[-1pt]
    \textbf{技术栈：}#3\\[-1pt]
    \textbf{项目链接：}\href{#4}{\nolinkurl{#4}}%
}

\begin{document}

%==================== 个人信息：三栏布局 ====================

\noindent
\begin{minipage}[c]{0.47\textwidth}
    {\LARGE\textbf{孔令航}}\\[5pt]
    电话：15832895339\\[2pt]
    邮箱：
    \href{mailto:17736236765@163.com}
    {17736236765@163.com}\\[2pt]
    GitHub：
    \href{https://github.com/Kong-lh-rgb}
    {github.com/Kong-lh-rgb}
\end{minipage}
\hfill
\begin{minipage}[c]{0.31\textwidth}

    \small
    \textbf{求职意向}\\[2pt]
    Agent 开发 / 算法实习生\\[7pt]
    \textbf{可实习时间}\\[2pt]
    \textbf{6 个月及以上，两周内到岗}
\end{minipage}
\hfill
\begin{minipage}[c]{0.14\textwidth}
    \raggedleft
    \includegraphics[
        width=2.30cm,
        height=3.22cm,
        keepaspectratio
    ]{resume-photo-blue.png}
\end{minipage}

\vspace{-10pt}

%==================== 教育背景 ====================

\section{教育背景}

\begin{tabularx}{\textwidth}{@{}X X c r@{}}
    \textbf{内蒙古大学（211）}
    & 计算机技术
    & 硕士
    & 2025--2028\\
    \textbf{河北金融学院}
    & 计算机科学与技术
    & 本科
    & 2021--2025\\
\end{tabularx}

%==================== 项目经历 ====================

\section{项目经历}

\projectTitle
{Vesta｜面向长程任务的本地 AI Agent Harness}
{2026.07--至今}
{Python、FastAPI、SQLite、MCP、Electron、React、TypeScript、Swift、OpenAI/Anthropic SDK}
{https://github.com/Kong-lh-rgb/vesta}

\vspace{2pt}

独立从 0 到 1 设计并实现本地 AI Agent Harness，
打通多模型适配、Agent Loop、Tool Calling、长程任务、上下文治理、
长期记忆、安全沙箱、Checkpoint/Recovery 与 Desktop 完整链路。

\begin{itemize}

\resumeItem{
    \textbf{Agent Runtime 与工具执行：}
    抽象统一 Model Adapter、Agent Loop 与
    \textbf{Registry--Executor--Hook} 工具执行层，
    编排本地工具、MCP、Skill 与 Computer Use；
    建立参数校验、超时、审批和沙箱边界，
    通过 WebSocket + JSON-RPC 打通 Electron、Python Host 与 Swift Helper。
}

\resumeItem{
    \textbf{长程任务、上下文与故障恢复：}
    设计 \textbf{Conversation--Run--Task} 分层状态模型，
    基于 Checkpoint 与 Recovery Run 支持模型异常、用户中断及 Host 重启后的任务恢复；
    通过分层裁剪、滚动压缩与 Run 级 Token Budget，
    将平均可计费 Token 从 \textbf{3,766 降至 2,767（-26.5\%）}，
    Prompt Cache 命中率由 \textbf{68.0\% 提升至 75.5\%}。
}

\resumeItem{
    \textbf{长期记忆与经验学习：}
    构建 \textbf{Core/Ordinary Memory} 分层记忆，
    结合 FTS5 与向量语义检索并通过 RRF 融合排序，
    以 Post-Run Reflection 和 Revision 实现跨会话更新；
    从 Completed Task 与真实 Trace 中提炼候选 Skill，
    经 Human Gate 审核后沉淀为可复用经验。
}

\resumeItem{
    \textbf{Agent Eval 与可靠性：}
    构建端到端 \textbf{Eval Harness}，记录 Trace、中间状态、失败原因与版本化 Baseline，
    建立“评测--定位--优化--回归”闭环；
    针对 68 个能力单元各执行 3 次，共完成 \textbf{204 次真实模型回归}，
    达到 \textbf{192/204 样本通过、57/68 能力单元连续三次通过}。
}

\end{itemize}

\vspace{-5pt}
\noindent\textcolor{muted}{\rule{\textwidth}{0.4pt}}
\vspace{-5pt}

\projectTitle
{Research Copilot｜金融投研多智能体系统（全栈）}
{2026.04--2026.06}
{Python、LangGraph、FastAPI、PostgreSQL、MCP、SSE、Next.js、Docker}
{https://github.com/Kong-lh-rgb/research-copilot}

\vspace{2pt}

构建面向 A 股研究的多智能体系统，
将自然语言研究问题拆解为可追踪的任务 DAG，并生成结构化研报。

\begin{itemize}

\resumeItem{
    \textbf{多智能体编排引擎：}
    基于 LangGraph 搭建
    \textbf{Controller--Planner--Worker--Reviewer} 状态图，
    通过 Send API 与依赖状态驱动任务并发执行和下游节点解锁，
    完成意图路由、规划、执行与结果审查链路。
}

\resumeItem{
    \textbf{工具接入与全栈交付：}
    基于 MCP 统一封装投研工具，持久化工作流状态并为敏感操作加入
    Human-in-the-Loop 审批；
    打通 FastAPI、SSE 与 Next.js 实时任务界面，
    使用 Docker Compose 完成 API、前端与 PostgreSQL 一键部署。
}

\end{itemize}

%==================== 开源贡献 ====================

\section{开源贡献}

\begin{tabularx}{\textwidth}{@{}X>{\raggedleft\arraybackslash}p{0.18\textwidth}@{}}
    \textbf{\large
        \href{https://github.com/multica-ai/multica}
        {Multica｜开源 AI Agent 协作平台（48.8k+ Stars）}}
    &
    \textcolor{muted}{2026.09}
\end{tabularx}

\vspace{1pt}

\textbf{核心代码贡献：}
修复共享浅克隆缓存遗漏最新提交、导致 Agent 隔离工作目录检出失败的问题；
调整 Repo Cache 引用同步与 Checkout 顺序并构造确定性回归测试，
完整后端 CI 通过，\href{https://github.com/multica-ai/multica/pull/8024}{PR \#8024} 已由维护者合并。

%==================== 专业技能 ====================

\section{专业技能}

\skill{编程与后端工程}{
    Python、Go、SQL；FastAPI、SQLite、PostgreSQL、Docker、Linux；
    熟悉异步编程、并发控制、SSE/WebSocket 与常见后端服务开发
}

\skill{LLM 与 Agent 系统}{
    OpenAI/Anthropic SDK、Structured Output、Function/Tool Calling、ReAct、MCP、LangGraph；
    Agent Runtime/Harness、上下文治理、长期记忆、Skill Learning、
    Human-in-the-Loop、Checkpoint/Recovery
}

\skill{Agent 工程与可靠性}{
    多模型接入、工具注册与执行、安全沙箱、超时/重试、故障恢复；
    Trace/日志分析、Agent Eval、Offline/Live Regression、
    LLM-as-a-Judge、规则与状态断言、Token/缓存分析、GitHub Actions
}

\skill{模型与知识检索}{
    Transformer、PyTorch、Hugging Face Transformers、LoRA/SFT；
    RAG、Embedding、BGE-M3、Milvus、Hybrid Search、Rerank、KG-RAG
}

\end{document}
